\documentclass[12pt,a4paper,twoside]{article}

% Common preamble I use
\usepackage[utf8]{inputenc}         % Input encoding
\usepackage[T1]{fontenc}            % Font encoding
\usepackage{textcomp}               % Provides various additional symbols and text-related features
\usepackage{microtype}              % Improves micro-typography of the document
\usepackage{pifont}                 % Provides additional symbols, including dingbats
\usepackage{changepage}             % Allows adjustment of page layout parameters
\usepackage{ftnxtra}                % Allows for fancy footnotes
\usepackage{fancyhdr}               % Allows for headers and footers
\setlength{\headheight}{14.49998pt}
\addtolength{\topmargin}{-2.49998pt}
\usepackage{subcaption}             % Enhanced support for sub-figures and sub-captions
\usepackage{wrapfig}

\usepackage{chemgreek,textgreek}    % Greek letters 

\usepackage{authblk}                % For customizing author block layout in the title page

% Listings setup
\usepackage{listings}               % Code listings
\usepackage{xcolor}                 % Color support
% Define R language for listings
\lstdefinelanguage{R}{
    keywords={function,
              if,
              else,
              repeat,
              while,
              for,
              in,
              next,
              break,
              TRUE,
              FALSE,
              NULL,
              Inf,
              NaN,
              NA,
              NA_integer_,
              NA_real_,
              NA_complex_,
              NA_character_,
              library,
              require,
              install.packages},
    keywordstyle=\color{blue}\bfseries,
    sensitive=true,
    comment=[l]{\#},
    commentstyle=\color{green!50!black},
    stringstyle=\color{orange},
    morestring=[b]'',
    morestring=[b]',
    moredelim=[is][\color{purple}]{`}{`}
}
\lstdefinestyle{Rstyle}{
    language=R,                
    basicstyle=\ttfamily\small,
    keywordstyle=\color{blue}, 
    stringstyle=\color{orange},
    commentstyle=\color{green!50!black}, 
    backgroundcolor=\color{gray!10},    
    showstringspaces=false,    
    numbers=left,              
    numberstyle=\tiny,         
    breaklines=true,           
    frame=single,              
}

% Graphic packages
\usepackage{graphicx}               % Graphics support
\usepackage[a4paper, margin=1in]{geometry}
\usepackage{tikz}                   % For in-doc drawings
\usepackage{pgfplots}
\pgfplotsset{compat=1.18}

% Maths packages
\usepackage{amsmath}                % Math support
\usepackage{amssymb}                % Provides various additional mathematical symbols
\usepackage{amsthm}                 % Provides enhanced support for theorem-like environments

% Table support packages
\usepackage{xtab}                   % Tables with adjustable width 
\usepackage{multirow}               % For multi-row cells in tables
\usepackage{array}                  % Provides more flexible and customizable array and tabular environments
\usepackage{longtable}              % Allows tables that span multiple pages
\usepackage{threeparttablex}        % Helps avoid clashes between tables and footnotes
\usepackage{tabularray}             % Combines features of longtable, array, multirow, etc.

\usepackage{booktabs}               % Enhanced tables
\usepackage{csquotes}               % Quotation marks and quote formatting
\usepackage{siunitx}                % For typesetting numbers and units.
\sisetup{                           % Centers numbers in tables by the decimal point.
  input-ignore={,},
  input-decimal-markers={.},
  group-separator={,},
  table-align-text-pre=false,
  table-align-text-post=true
  }

\usepackage{algorithm}              % Allows insertion of algorithms as an object
\usepackage{algpseudocode}          % Used to insert pseudocode  
\usepackage{caption}                % Customizes captions in floating environments

\usepackage{setspace}

\usepackage{ragged2e}               % Enhanced text alignment commands

\usepackage{enumitem}               % Customizable lists     

\usepackage[sort&compress]{natbib}  % Cite style
\setcitestyle{numbers,square,comma}

\usepackage{times}                  % Times font       

\usepackage{hyperref}               % Hyperlinks
\usepackage{url}                    % For typesetting URLs with line breaks at hyphens
\def\UrlBreaks{\do\/\do-}           % Allows urls to break lines chktex 4
\hypersetup{breaklinks=true}
\urlstyle{same}
\usepackage{nameref}                % Enables referencing section names instead of numbers

% For using TODO notes
% \todo[inline,caption={}]{TODOs are to be inserted like this}
%\usepackage[color=blue!10,textsize=footnotesize,textwidth=25mm]{todonotes}
%\usepackage[disable]{todonotes}
\usepackage{float}

% CUSTOM LABELS
\DeclareCaptionLabelFormat{apafigure}{\textbf{Figure #2}}
\DeclareCaptionLabelFormat{procedure}{\textbf{Procedure #2}}
\DeclareCaptionLabelFormat{apalisting}{\textbf{Listing #2}}

% CUSTOM CAPTION FORMATS
\captionsetup[figure]{labelformat=apafigure,
                      labelsep=newline,
                      justification=justified,
                      singlelinecheck=false}
\DeclareMathOperator{\atan}{atan}
\captionsetup[lstlisting]{labelformat=apalisting,
                          position=above,
                          justification=raggedright,
                          labelsep=newline,
                          singlelinecheck=false,}
\DeclareCaptionLabelSeparator*{spaced}{\\[2ex]}
\captionsetup[table]{textfont=bf,
                     format=plain,
                     justification=justified,
                     singlelinecheck=false,
                     labelsep=newline,
                     skip=0pt}

\makeatletter
% Commands to format a counter value as Greek letter to be used like 
% \arabic or \roman when using enumerate:
\newcommand*\alphgreek[1]{\expandafter\@alphgreek\csname c@#1\endcsname}
\newcommand*\@alphgreek[1]{\csname chemgreekIntToGreek:n\endcsname{#1}}
\newcommand*\Alphgreek[1]{\expandafter\@Alphgreek\csname c@#1\endcsname}
\newcommand*\@Alphgreek[1]{\csname chemgreekIntToGreek:n\endcsname{#1}}

% Register new counter formats to enumitem:
\AddEnumerateCounter*{\alphgreek}{\@alphgreek}{\chemalpha}
\AddEnumerateCounter*{\Alphgreek}{\@Alphgreek}{\chemAlpha}
\makeatother

% \code{} command to insert in-text code snippets
\definecolor{light-gray}{gray}{0.95}
\newcommand{\code}[1]{\colorbox{light-gray}{\texttt{#1}}}

% AUTHOR
\newcommand{\authorname}{Samuel Jonsson}
\newcommand{\authormail}{\href{mailto:sajs19@student.bth.se}{sajs19@student.bth.se}}
\newcommand{\authorpersonnr}{19990415-5596} % chktex 8
\newcommand{\program}{DVAMI19h}

% REPORT METADATA
\newcommand{\doctitle}{MS2505: Bayesian Statistics}
\newcommand{\docsubtitle}{Course Project}
\newcommand{\docdate}{Latest version: \today}

\pagestyle{fancy}                                 % Use the fancy page style
\fancyhf{}                                        % Clear all headers and footers

% HEADER CUSTOMISATION
\fancyhead[L]{\textit{\doctitle~--~\docsubtitle}} % Header float left
\fancyhead[R]{\textit{\docdate}}                  % Header float right

% FOOTER CUSTOMISATION
\fancyfoot[L]{\authorname~--~\authormail}         % Footer float left
\fancyfoot[R]{Page \thepage}                      % Footer float right

\setlist{topsep=1ex,itemsep=0.5ex,parsep=0pt,partopsep=0pt}

\begin{document}

\vspace*{3cm}

\begin{center}
    \Huge\textbf{\doctitle}\\
    \Large\textbf{\docsubtitle}
    \\\vspace*{5mm} 
    \large\today                            
\end{center} 

\vspace*{5mm}

\begin{tabular}{|l|p{13cm}|}
    \hline
    \textbf{Name}           & \authorname     \\
    \hline
    \textbf{E-Mail}         & \authormail     \\
    \hline
    \textbf{Person Nr.}     & \authorpersonnr \\
    \hline
    \textbf{Program}        & \program        \\
    \hline
\end{tabular}

\begin{figure}[!b]
    \centering
    \includegraphics[width = 0.25\textwidth]{figures/BTH_logo_black.png}
\end{figure}

\newpage

\section{Setup}
\label{sec:setup}
\todo[inline,caption={}]{
    \begin{itemize}
        \item Describe the data and the analysis problem.
        \item Choose and describe the modeling approach (e.g., non-hierarchical or hierarchical model).
        \item Justify your prior choice.
        \item Perform posterior predictive checks.
    \end{itemize}
}
\subsection{Analysis problem}
\label{ssec:analysisproblem}

\subsection{Data Selection}
\label{ssec:dataselection}
\todo[inline,caption={}]{Describe the data and the analysis problem.}

The dataset selected is a datasets containing a list of emails, as well as a label marking each
email as "spam" or "ham" (spam or not spam). The first 10 rows of the dataset looks as follows:

\begin{longtable}[!ht]{p{1.5cm}|p{11cm}}
    \caption{\code{mail\_data.csv} dataset first 10 rows}\\
    \textbf{Category} & \textbf{Message} \\
    \hline
    ham      & "Go until jurong point, crazy.. Available only in bugis n great world la e buffet...
               Cine there got amore wat..." \\
    \hline
    ham      & Ok lar... Joking wif u oni... \\
    \hline
    spam     & Free entry in 2 a wkly comp to win FA Cup final tkts 21st May 2005. Text FA to 87121 to
               receive entry question(std txt rate)T\&C's apply 08452810075over18's \\
    \hline
    ham      & U dun say so early hor... U c already then say... \\
    \hline
    ham      & "Nah I don't think he goes to usf, he lives around here though" \\
    \hline
    spam     & "FreeMsg Hey there darling it's been 3 week's now and no word back! I'd like some fun you
                up for it still? Tb ok! XxX std chgs to send, \pounds 1.50 to rcv" \\
    \hline
    ham      & Even my brother is not like to speak with me. They treat me like aids patent. \\
    \hline
    ham      & As per your request 'Melle Melle (Oru Minnaminunginte Nurungu Vettam)' has
               been set as your callertune for all Callers. Press *9 to copy your friends Callertune \\
    \hline
    ...      & ... \\
\end{longtable}

Then, using a Python script, the labels were converted to 1 if it was "spam" and 0 if it was "ham",
for easier analysis.

\subsection{Model}
\label{ssec:model}
\todo[inline,caption={}]{
    \begin{itemize}
        \item Choose and describe the modeling approach (e.g., non-hierarchical or hierarchical model).
        \item Justify your prior choice.
    \end{itemize}
}
The model chosen was a binomial likelihood model with a beta prior. As the goal is to analyse the
probability of an email being spam, the fallout will be binary (either it is spam or it is not). Hence,
a binomial likelihood, where I want to find the parameter $\theta$ in a dataset of fixed size with a set
number of "successes" and "fails" (spam and ham), is appropriate.

Additionally, as I do not have any prior knowledge in regards to this distribution, a non-informative prior
is the most suited option, and as $Beta(1,1)$ is a common prior used with binomial likelihood functions, I
chose it for this problem.

\subsection{Prior checks}
\label{ssec:priorchecks}
\todo[inline,captio?{}]{Perform posterior predictive checks.}

\section{Results}
\todo[inline,caption={}]{Include diagnostics to assess model convergence and adequacy.}

\section{Discussion}
\todo[inline,caption={}]{Discuss results, problems encountered, and possible improvements.}

\appendix
\section{R Code}
\lstinputlisting[language=R,style=Rstyle,caption=Project R code,label=apx:RCode,]{./R/base.r}
\end{document}